\section{Introduction}

Speech emotion recognition (SER) aims to identify emotional states from acoustic signals. While substantial progress has been made on adult speech corpora using self-supervised learning (SSL) models, children's speech emotion recognition remains underexplored despite its importance for educational technology, clinical screening, and child-robot interaction.

Children's speech differs from adult speech along multiple acoustic dimensions---higher and more variable fundamental frequency (F0), shorter utterance durations, and ongoing anatomical development of the vocal tract. These differences raise a fundamental question: is the performance gap between children's and adult SER driven by model architecture choices, or by inherent distributional properties of the speech signal itself?

Prior work has approached children's SER through two main routes: applying SSL models pre-trained on adult speech directly to children's data, and designing specialized pooling or adaptation modules. However, systematic diagnosis of \textit{why} performance varies across children's corpora has been lacking. In particular, the contribution of distribution shift---the statistical mismatch between training and testing distributions---has not been isolated from architectural factors in a controlled experimental framework.

In this work, we present a systematic 192-experiment ablation study that disentangles the effects of distribution shift from architectural design choices in children's SER. Using WavLM Base as a unified backbone across three corpora---C-BESD (children's acted speech), FAU Aibo (children's spontaneous speech), and IEMOCAP (adult acted speech)---we evaluate pooling strategies, layer fusion, adapter modules, data augmentation, and transfer learning under both frozen and unfrozen conditions.

Our key findings are: (1) distribution shift between corpora, not architecture, is the dominant bottleneck---zero-shot cross-corpus performance collapses to 19--35\% even with optimal pooling; (2) the acted-versus-spontaneous distinction within children's speech creates a larger performance gap (25pp) than the child-versus-adult distinction; (3) transfer learning success is determined primarily by the target domain, with C-BESD as a target achieving 91\%+ regardless of source; (4) adapter modules provide no benefit and can be safely removed; and (5) layer fusion contributes negligibly, as any single transformer layer L$\geq$7 performs equivalently to learned weighted fusion.

We release all 192 experiment logs, model checkpoints, and evaluation code to facilitate reproducible research on children's SER.
